\documentclass{article}


\PassOptionsToPackage{numbers,compress}{natbib}


% ready for submission
\usepackage[preprint]{neurips_2026}


\usepackage[utf8]{inputenc} % allow utf-8 input
\usepackage[T1]{fontenc}    % use 8-bit T1 fonts
\usepackage{hyperref}       % hyperlinks
\usepackage{url}            % simple URL typesetting
\usepackage{booktabs}       % professional-quality tables
\usepackage{amsfonts}       % blackboard math symbols
\usepackage{amsmath}
\usepackage{nicefrac}       % compact symbols for 1/2, etc.
\usepackage{microtype}      % microtypography
\usepackage{xcolor}         % colors
\newcommand{\new}[1]{{\color{red} #1}}


\title{ECE228 Project Proposal:\\
Refining Video Object Segmentation with Test-time Gradient Correction}


\author{%
  Suyu Chen \\
  \texttt{suc033@ucsd.edu}
  \And
  Yi Sun \\
  \texttt{yis094@ucsd.edu}
  \And
  Yinru Chen\\
  \texttt{yic173@ucsd.edu}
  \AND
  \textbf{Group \#71}
}


\begin{document}


\maketitle


\section{Problem Background \& Motivation}

Semi-supervised video object segmentation (VOS) models~\citep{pont2017davis}, such as STM~\citep{oh2019stm} and AOT~\citep{yang2021aot}, predict target masks across video sequences given a first-frame annotation. Currently, state-of-the-art pipelines are evaluated predominantly by per-frame spatial overlap ($\mathcal{J}\&\mathcal{F}$). This metric, however, is blind to temporal consistency: a temporally jittery mask scores similarly to a smooth one. In physical robotics applications like visual servoing~\citep{chaumette2006visual}, this high-frequency perception noise is disastrous, translating directly into physical actuator jerk, wasted control energy, and tracking lock-loss.

To bridge the gap between perception quality and physical control stability, we optimize VOS models via \emph{test-time gradient correction}~\citep{li2020cycle}. At inference time, this technique uses the first-frame ground truth as an anchor, iteratively refining the current frame's prediction through gradient descent on a cycle-consistency loss back to frame $1$.

In this project, we apply test-time gradient correction to optimize the AOT~\citep{yang2021aot} pipeline. Crucially, we evaluate the system not just on standard VOS metrics, but within a simulated closed-loop visual-servoing environment to determine if this optimization yields measurable stability improvements in physical tracking tasks.

\section{Related Work}

\paragraph{Memory-based VOS architectures.}
The dominant family for semi-supervised VOS uses an external memory of
past frame--mask pairs and attention-based matching.
STM~\citep{oh2019stm} introduced this paradigm; AOT~\citep{yang2021aot}
added an identification mechanism that processes multiple objects in a
single forward pass; DeAOT~\citep{yang2022deaot} decoupled visual and
ID features for hierarchical propagation, and
XMem~\citep{cheng2022xmem} added a long-term memory pool for
minute-long videos.  All four operate purely \emph{offline} at
inference, leaving room for test-time refinement.

\paragraph{Test-time refinement and cycle consistency in VOS.}
Online-learning approaches such as OnAVOS~\citep{voigtlaender2017onavos}
fine-tune network weights on the first frame; they are accurate but
prohibitively slow ($<\!1$ FPS).  \citet{li2020cycle} instead use the
first-frame ground-truth mask as a self-supervisory anchor and apply
\emph{gradient correction}: a few PGD-style updates of the predicted
mask along the gradient of a cycle-consistency loss
$\mathcal{L}\bigl(S_\theta(X_t,\hat Y^{l}_t,X_1),Y_1\bigr)$.  Cycle
consistency has a long history in image translation, but
\citet{li2020cycle} is to our knowledge the only work that uses it as a
test-time \emph{control signal} for VOS.  Our project sits in this
thread but re-evaluates the technique through a control-theoretic lens.

\paragraph{Perception for visual servoing.}
Classical image-based visual servoing (IBVS)~\citep{chaumette2006visual}
reduces tracking control to driving a per-frame feature error to zero.
Modern systems typically plug in a bounding-box deep tracker, not
pixel-precise VOS, and seldom analyse how perception errors propagate
into control errors.  We use IBVS as a downstream evaluation harness
for VOS, which has not been done in the gradient-correction literature.

\section{High-level Methodology}

\paragraph{Step 1 -- Reproduce and audit.}
We reproduce the public AOT $+$ gradient-correction pipeline on the
DAVIS-2017 validation split and audit it against \citet{li2020cycle}
Eq.~(5), resulting J and F scores may act as indicators of the correctness of our implementation.

\paragraph{Step 2 -- Improvements.}
We then explore improvements on the corrected baseline:

\textbf{(I-1)} restore an end-to-end differentiable cycle by removing
the detach calls and using a straight-through estimator through the
$\arg\max$;

\textbf{(I-2)} introduce a \emph{multi-anchor} cycle that,
in addition to $Y_1$, uses a small set of high-confidence intermediate
predictions as auxiliary anchors to reduce appearance drift on long
sequences.

\paragraph{Step 3 -- Visual-servoing evaluation.}
Each predicted mask sequence is fed into a closed-loop simulator that
extracts the per-object centroid, smooths it through a 2-D
constant-velocity Kalman filter, and uses a discrete PID controller to
drive a virtual pan--tilt camera toward the centroid.

Because the harness's discriminative power depends on scene dynamics,
we stratify across four DAVIS buckets --- single-target with
distractors (\texttt{car-roundabout}), multi-target
(\texttt{bike-packing}, \texttt{cows}), high-dynamics
(\texttt{motocross-jump}) and long-drift (\texttt{breakdance}). 

\paragraph{Risks.}
If controller-level deltas turn out small once the implementation is
fixed, the project still delivers a controlled negative result, an
audited reference implementation, and an open-source evaluation
harness.

\bibliographystyle{plainnat}
{\footnotesize
\setlength{\bibsep}{0pt plus 0.2ex}
\begin{thebibliography}{9}
\providecommand{\natexlab}[1]{#1}

\bibitem[Chaumette and Hutchinson(2006)]{chaumette2006visual}
F.~Chaumette and S.~Hutchinson.
\newblock Visual servo control. {I}. {B}asic approaches.
\newblock \emph{IEEE Robotics \& Automation Magazine}, 13(4):82--90, 2006.

\bibitem[Cheng and Schwing(2022)]{cheng2022xmem}
H.~K. Cheng and A.~G. Schwing.
\newblock {XM}em: Long-term video object segmentation with an {A}tkinson--{S}hiffrin memory model.
\newblock In \emph{ECCV}, 2022.

\bibitem[Li et~al.(2020)Li, Xu, Peng, See, and Lin]{li2020cycle}
Y.~Li, N.~Xu, J.~Peng, J.~See, and W.~Lin.
\newblock Delving into the cyclic mechanism in semi-supervised video object segmentation.
\newblock In \emph{NeurIPS}, 2020.

\bibitem[Oh et~al.(2019)Oh, Lee, Xu, and Kim]{oh2019stm}
S.~W. Oh, J.-Y. Lee, N.~Xu, and S.~J. Kim.
\newblock Video object segmentation using space-time memory networks.
\newblock In \emph{ICCV}, 2019.

\bibitem[Pont-Tuset et~al.(2017)]{pont2017davis}
J.~Pont-Tuset, F.~Perazzi, S.~Caelles, P.~Arbel\'aez, A.~Sorkine-Hornung, and L.~Van~Gool.
\newblock The 2017 {DAVIS} challenge on video object segmentation.
\newblock \emph{arXiv:1704.00675}, 2017.

\bibitem[Voigtlaender and Leibe(2017)]{voigtlaender2017onavos}
P.~Voigtlaender and B.~Leibe.
\newblock Online adaptation of convolutional neural networks for video object segmentation.
\newblock In \emph{BMVC}, 2017.

\bibitem[Yang et~al.(2021)Yang, Wei, and Yang]{yang2021aot}
Z.~Yang, Y.~Wei, and Y.~Yang.
\newblock Associating objects with transformers for video object segmentation.
\newblock In \emph{NeurIPS}, 2021.

\bibitem[Yang and Yang(2022)]{yang2022deaot}
Z.~Yang and Y.~Yang.
\newblock Decoupling features in hierarchical propagation for video object segmentation.
\newblock In \emph{NeurIPS}, 2022.

\end{thebibliography}
}


\end{document}
